\chapter{Conclusions and Future Works}

This internship at Wells Fargo provided practical exposure to enterprise-scale engineering where reliability, governance, and response speed are as important as feature delivery. Through the three workstreams, I contributed to decisioning modernization, cost-aware architecture strategy, and fraud data centralization.

\section{Conclusions}

The major conclusions are as follows: automating rule compilation and artifact publication substantially improves consistency and reduces manual operational risk, repository consolidation improves maintainability and lowers cognitive overhead for development teams, in-house stand-in decisioning capability is a viable strategic direction for reducing long-term vendor dependency and improving service continuity, and event-driven fraud centralization is essential for maintaining a unified risk view across teams and preventing contradictory decisions. Overall, the internship demonstrated that sustainable impact comes from systems thinking: aligning architecture choices with business goals, compliance constraints, and operational realities.

\section{Future Works}

Future scope for this work includes:

\begin{enumerate}
	\item \textbf{Decisioning automation hardening}
	\begin{itemize}
		\item Add rule schema validation and policy quality gates before compilation.
		\item Introduce approval-driven promotion flow across environments.
	\end{itemize}

	\item \textbf{Stand-in server productionization}
	\begin{itemize}
		\item Build parity dashboards comparing internal and third-party outcomes.
		\item Execute phased traffic migration with rollback guardrails.
	\end{itemize}

	\item \textbf{Fraud pipeline resilience}
	\begin{itemize}
		\item Strengthen idempotency controls and event contract governance.
		\item Add dead-letter queue analytics and lag-based alerting.
	\end{itemize}
\end{enumerate}

\section{Personal Learning Outcomes}

\subsection{Technical and Professional Learning}

I learned to translate business pain points into implementable technical architecture, design migration paths with controlled risk, and evaluate solutions through enterprise non-functional requirements such as reliability, traceability, and scalability.

\subsection{Learning to Use GitHub Copilot Better}

During this internship, I improved how I use GitHub Copilot in practical engineering tasks. I learned to write context-rich prompts so generated code better matches domain requirements, used Copilot to accelerate boilerplate, refactoring, and documentation drafting while keeping validation responsibility with myself, employed iterative prompting to compare alternatives for readability, maintainability, and production safety, and learned that Copilot is most effective as a productivity partner when paired with strong engineering judgment and testing discipline. These learnings improved both development speed and communication quality during project execution.